\documentclass[portrait]{seminar}
\usepackage{colordvi,boldmath,semcolor,semlayer,rotate,epsf,fancybox,mathsym}

\slidewidth9.5in
\slideheight6.7in
%\def\slidetopmargin{\iflandscape.2in\else0.8in\fi}
\def\slidetopmargin{\iflandscape.3in\else0.8in\fi}
\def\slidebottommargin{\iflandscape1.0in\else0.4in\fi}
\def\slideleftmargin{\iflandscape.2in\else.5in\fi}
\def\sliderightmargin{\iflandscape1.0in\else.8in\fi}

\newcommand{\ket}[1]{|#1\rangle}
\newcommand{\bra}[1]{\langle #1|}
\newcommand{\scalar}[2]{\langle #1|#2\rangle}

\newcommand{\cH}{{\cal H}}

\slidesmag{4}
\centerslidesfalse
\slideframewidth0.5pt
%\slideframe{none}

%\portraitonly
%\landscapeonly
%\def\leftsliderotation#1{\rotate[l]{#1}}
%\sliderotation{left}
\rotateheaderstrue


\newcommand{\onemat}{{\bf 1}}

\pagestyle{empty}

\begin{document}

\begin{slide}
{$ $}
\vfill
\begin{center}
{\large \blue Estimating mixing properties of local Hamiltonian
dynamics and continuous quantum random walks is PSPACE-hard}

\vfill {\large \blue Measuring local observables could probabilistically
solve PSPACE}

\vfill Pawel Wocjan

{\small Institut f\"ur Algorithmen und Kognitive Systeme}

{\small Universit\"at Karlsruhe}
\end{center}
\vfill
\end{slide}

%%%%%%%%%%%%%%%%%%%%%%%%%%%%%%%%%%%%
%
% 
%
%%%%%%%%%%%%%%%%%%%%%%%%%%%%%%%%%%%%

\begin{slide}
{\large {\bf Observable, Hamiltonian, spectral resolution}}

\vspace{0.25cm}
finite dimensional Hilbert space ${\cal{H}}$

\vspace{0.25cm}
{\blue {\bf observable}}\, /\, {\blue {\bf Hamiltonian}}\, :
\, Hermitian operator $A$ on ${\cal{H}}$

{\blue {\bf continuous quantum random walk}}\, : $A$ is an adjacency
matrix

\vspace{0.25cm}
spectral resolution of $A$:

\begin{itemize}
\item different eigenvalues of $A$: $(\lambda)$ 
\item orthogonal projections onto the corresponding eigenspaces of $A$:
$(P_{\lambda})$
\end{itemize}
\end{slide}

%%%%%%%%%%%%%%%%%%%%%%%%%%%%%%%%%%%%
%
% 
%
%%%%%%%%%%%%%%%%%%%%%%%%%%%%%%%%%%%%

\begin{slide}
{\large {\bf Mixing properties of Hamiltonian dynamics and continuous
quantum random walks}}

\vspace{0.5cm} 
{\blue {\bf Time-average of a state $\rho$ on $\cH$}}\,:
\[
\bar{\rho}:=\lim_{T\rightarrow\infty} 
\frac{1}{T} \int_{t=0}^T e^{-i A t}\,\rho\, e^{i A t}\, {\rm d}t
\]

{\blue {\bf von-Neumann entropy $S(\bar{\rho})$ of the time-average}}
\[
S(\bar{\rho}):= {\rm tr} (\bar{\rho}\log_2\bar{\rho})
\]

motivation for considering $S(\bar{\rho})$:
\begin{itemize}
\item quantifies ``how much and how uniformly'' the orbit $\rho_t:=e^{-i
A t}\,\rho\, e^{i A t}$ fills out the state space

\item quantifies ``how far'' a pure state $\rho:=\ket{\Phi}\bra{\Psi}$
is from all eigenvectors
\end{itemize}
\end{slide}

\begin{slide}
{\large {\bf Measurement of observables}}

\vspace{0.5cm} 
{\blue {\bf Postulate: Measurement of $A$ with accuracy $\Delta$}}\,:

for all states $\rho$ on $\cH$ the probability to obtain an outcome in
the interval
\[
[\lambda-\Delta,\lambda+\Delta]
\]
is a least 
\[
(3/4)\, {\rm tr}(\rho P_{\lambda})
\]
\end{slide}

\begin{slide}
{\large {\bf Local operators}}

\vspace{0.5cm}
$\cH$ tensor product space of $n$ Hilbert spaces
\[
\cH:=\cH_1\otimes\cH_2\otimes\cdots\otimes\cH_n
\]

\vspace{0.5cm}
{\blue {\bf $A$ is $k$-local}}\,:

$A$ sum of operators acting on at most $k$ tensor components
non-trivially

\vspace{0.5cm}
example for $2$-local Hamiltonians:

Ising interactions, coupling topology specified by a graph $G=(V,E)$
\[
\sum_{(k,l)\in E} \sigma_z^{(k)}\sigma_z^{(l)}
\]
$\sigma_z^{(k)}=\onemat\otimes\cdots\otimes\onemat\otimes
\underbrace{\quad\sigma_z\quad}_{\mbox{$k$th qubit}}
\otimes\onemat\otimes\cdots\otimes\onemat$
\end{slide}

\begin{slide}
{\bf we want to prove}

\vfill

\begin{center}
{\large \blue Estimating mixing properties of local Hamiltonian
dynamics and continuous quantum random walks is PSPACE-hard}

\vfill

{\large \blue Measuring local observables could probabilistically
solve PSPACE}

\vfill
\end{center}

Outline of the proof:
\begin{itemize}
\item recall definition of PSPACE (based on irreversible Turing machines)
\item transform irreversible TMs to reversible TM
\item construct quantum circuit that simulate the computational steps
of the reversible TMs
\item encode the sequence of gates of these quantum circuits in local
operators

$\longrightarrow$ desired observables, Hamiltonians, adjacency
matrices
\end{itemize}
\vfill
\end{slide}

\begin{slide}
{\bf {\large Definition of PSPACE}}

class of all languages recognizable by {\blue polynomial space bounded} DTMs

Deterministic Turing Machine:
\begin{itemize}
\item head (internal states $Q$) 

tape divided into cells (contain symbols of the alphabet $\Sigma$)
\item local transition function $Q\times \Sigma \rightarrow Q \times
\Sigma \times \{-1,0,+1\}$

\begin{picture}(75,40)
\thicklines
\put( 0, 0){\framebox(15,12){}}
\put(15, 0){\framebox(15,12){}}
\put(30, 0){\framebox(15,12){$a$}}
\put(45, 0){\framebox(15,12){}}
\put(60, 0){\framebox(15,12){}}
\put(30,20){\makebox(15,15){$p$}}
\put(37.5,27.5){\oval(15,15)}
{\red \put(13,-2){\dashbox{0.5}(49,39){}}}
\end{picture}
$\quad{\Large \Longrightarrow}\quad$
\begin{picture}(100,35)
\thicklines
\put( 0, 0){\framebox(15,12){}}
\put(15, 0){\framebox(15,12){}}
\put(30, 0){\framebox(15,12){$b$}}
\put(45, 0){\framebox(15,12){}}
\put(60, 0){\framebox(15,12){}}
\put(45,20){\makebox(15,15){$q$}}
\put(52.5,27.5){\oval(15,15)}
\end{picture}
\item input: $x_1\, x_2\,\ldots x_l$ \hspace{2.15cm}
output: symbol in current cell when\\
${}$\hspace{6.25cm} RTM goes into a final state $q_f$

\begin{picture}(140,40)
\thicklines
\put( 0, 0){\framebox(15,12){$x_1$}}
\put(15, 0){\framebox(15,12){$x_2$}}
%
\put(37, 0){\makebox(15,12){$\cdots$}}
%
\put(60, 0){\framebox(15,12){$x_l$}}
\put(0,20){\makebox(15,15){$q_i$}}
\put(7.5,27.5){\oval(15,15)}
\put(75, 0){\framebox(15,12){$\#$}}
\put(90, 0){\framebox(15,12){$\#$}}
\put(110, 0){\makebox(15,12){$\cdots$}}
%
{\blue 
\put(-2,-2){\dashbox{0.5}(140,39){}}
}
\put(60,39){\makebox(15,15){{\small {\blue $poly(l)$}}}}
\end{picture}
\quad\quad
\begin{picture}(75,40)
\thicklines
\put(-2, 0){\makebox(15,12){$\cdots$}}
\put(15, 0){\framebox(15,12){}}
\put(30, 0){\framebox(15,12){$y$}}
\put(45, 0){\framebox(15,12){}}
\put(62, 0){\makebox(15,12){$\cdots$}}
\put(30,20){\makebox(15,15){$q_f$}}
\put(37.5,27.5){\oval(15,15)}
%{\red \put(13,-2){\dashbox{0.5}(49,39){}}}
\end{picture}

\end{itemize}
\end{slide}

\begin{slide}
{\bf {\large Quantified Boolean Formula Problem}}

\vspace{0.25cm}
is a {\blue PSPACE-complete} problem

\vspace{0.25cm}
quantified Boolean formula:
\[
Q_1 v_1\, Q_2 v_2 \cdots Q_n v_n\, B(v_1,v_2,\ldots,v_n)
\]
\quad $Q_i$: quantifiers (either $\forall$ or $\exists$)

\quad $B(v_1,v_2,\ldots,v_n)$: Boolean formula (without quantifiers)

\hspace{3.15cm} in the variables $v_1,v_2,\ldots,v_n$

\vspace{0.25cm}
{\bf quantified Boolean formula problem:}

determine if a quantified formula is true

\vspace{0.25cm} 
it can be solved within space $O(l^2)$ ($l$ input length) by an
irreversible TM
\end{slide}

\begin{slide}
{\bf {\large Reversible Turing machines}}

{\footnotesize Lange et al., ``Reversible space equals deterministic space'',
J.~Comp.~\& Syst.~Sci., 2000.}

{\blue irreversible TMs can be simulated space-efficiently by reversible TMs}

explicit construction of the RTM

\vspace{0.25cm}
each transition is
\begin{itemize}
\item either {\red moving} $p\rightarrow (q,\pm 1)$

\begin{picture}(75,35)
\thicklines
\put( 0, 0){\framebox(15,12){}}
\put(15, 0){\framebox(15,12){}}
\put(30, 0){\framebox(15,12){}}
\put(45, 0){\framebox(15,12){}}
\put(60, 0){\framebox(15,12){}}
\put(30,20){\makebox(15,12){$p$}}
\put(37.5,27.5){\oval(15,15)}
\end{picture}
$\quad{\Large \Longrightarrow}\quad$
\begin{picture}(100,35)
\thicklines
\put( 0, 0){\framebox(15,12){}}
\put(15, 0){\framebox(15,12){}}
\put(30, 0){\framebox(15,12){}}
\put(45, 0){\framebox(15,12){}}
\put(60, 0){\framebox(15,12){}}
\put(45,20){\makebox(15,15){$q$}}
\put(52.5,27.5){\oval(15,15)}
\end{picture}
\item or {\red read-and-write} $(p,a)\rightarrow (q,b)$

\begin{picture}(75,35)
\thicklines
\put( 0, 0){\framebox(15,12){}}
\put(15, 0){\framebox(15,12){}}
\put(30, 0){\framebox(15,12){$a$}}
\put(45, 0){\framebox(15,12){}}
\put(60, 0){\framebox(15,12){}}
\put(30,20){\makebox(15,12){$p$}}
\put(37.5,27.5){\oval(15,15)}
\end{picture}
$\quad{\Large \Longrightarrow}\quad$
\begin{picture}(75,35)
\thicklines
\put( 0, 0){\framebox(15,12){}}
\put(15, 0){\framebox(15,12){}}
\put(30, 0){\framebox(15,12){$b$}}
\put(45, 0){\framebox(15,12){}}
\put(60, 0){\framebox(15,12){}}
\put(30,20){\makebox(15,12){$q$}}
\put(37.5,27.5){\oval(15,15)}
\end{picture}
\end{itemize}
\end{slide}

\begin{slide}
{\bf {\large Quantum circuit simulating the RTM\quad (first idea)}}

registers: \texttt{head}\,,\,\,\texttt{tape\_index}\,,\,\,
\texttt{ACC}\,\,,\,\,\texttt{cell[i]}\,
$i=0,\ldots,N$
\vfill
\begin{itemize}
\item $U_{\rm moving}$ implements $p\rightarrow (q,\pm1)$
\[
\ket{p}_{\rm head}\otimes\ket{i}_{\rm tape\_index}
\rightarrow
\ket{q}_{\rm head}\otimes\ket{i\pm 1}_{\rm tape\_index}
\]
\vfill
\item
$U_{\rm r\&w}$ implements $(p,a)\rightarrow (q,b)$
\begin{itemize}
\item swap \texttt{ACC} and \texttt{cell[i]} (current cell) 
\[
\prod_{i=0}^{N-1} 
\Lambda_i\big(
{\rm swap}({\rm ACC},{\rm cell[i]})\big)
\quad\quad
\mbox{{\small $\Lambda_i$ applies swap iff \texttt{tape\_index}=i}}
\]
\item apply $W_{\rm r\&w}$
\[
\ket{p}_{\rm head}\otimes\ket{a}_{\rm ACC} \rightarrow
\ket{q}_{\rm head}\otimes\ket{b}_{\rm ACC}
\quad\quad\quad\quad\quad\quad\quad\quad\quad\quad\quad\quad\quad{}
\]
\item swap \texttt{ACC} and \texttt{cell[i]}
\end{itemize}
\end{itemize}
\end{slide}

\begin{slide}
\begin{picture}(50,230)
\thicklines
\put(  0,200){\framebox(50,20){{\rm tape\_index}}}
\put(  0,160){\framebox(50,20){{\rm head}}}
\put(  0,120){\framebox(50,20){{\rm ACC}}}
\put(  0, 80){\framebox(50,20){{\rm cell 0}}}
\put(  0, 40){\makebox(50,20){$\vdots$}} %
\put(  0,  0){\framebox(50,20){{\rm cell N-1}}}
%
\put( 70,155){\framebox(50,70){$U_{\rm moving}$}}
%
\put(150,210){\circle*{5}}
\put(150,210){\line(0,-1){100}}
%
\put(140,130){\line(1,-2){20}}
\put(140,90){\line(1,2){20}}
%
\put(200,210){\circle*{5}}
\put(200,210){\line(0,-1){140}}
%
\put(190,130){\line(1,-6){20}}
\put(190, 10){\line(1, 6){20}}
%
% quantum wires
%
\put( 60,210){\line(1,0){10}}  % tape index
\put(120,210){\line(1,0){180}}
\put( 60,170){\line(1,0){10}}  % head
\put(120,170){\line(1,0){120}}
\put(290,170){\line(1,0){10}}
\put( 60,130){\line(1,0){80}}  % ACC
\put(210,130){\line(1,0){30}}
\put(290,130){\line(1,0){10}}
\put( 60, 90){\line(1,0){80}}  % cell 0
\put(210, 90){\line(1,0){90}}
\put( 60, 10){\line(1,0){130}} % cell N-1
\put(210, 10){\line(1,0){90}}
%
% dashed lines
\multiput(160,130)(5,0){6}{\line(1,0){2.5}}
\multiput(160, 90)(5,0){10}{\line(1,0){2.5}}
%
% W_r&w
%
\put(240,110){\framebox(50,70){$W_{{\rm r\&w}}$}}
\end{picture}
\end{slide}

\begin{slide}
{\bf {\large Disadvantages of the quantum circuit}}
\begin{itemize}
\item not known exactly: how many times $U$ should be applied 

desired: running time $r_l$ depends only on input length $l$
\item $U$ produces garbage (e.g.\ does not reset ancillas)

desired: 
\[
{\blue 
V^{r_l}\ket{x}\otimes\ket{{\red y}}\otimes\ket{00\ldots 0}=
\ket{x}\otimes\ket{{\red f(x)\oplus y}}\otimes\ket{00\ldots 0}
}
\]
\hspace{1.5cm}$f(x)$ solution to the instance $x$
\end{itemize}

\vspace{0.25cm}
$\Longrightarrow$ second idea: augment $U$ and $U^\dagger$ with some
control logic

\medskip
\centerline{\epsfxsize0.8\textwidth\epsfbox{operationmode.eps}}

\end{slide}

\begin{slide*}
{\bf Final quantum circuit}

\bigskip
\centerline{\epsfxsize1.0\textwidth\epsfbox{schaltkreis_1bit_neu.eps}}
\end{slide*}

\begin{slide}
{\bf {\large Characterization of PSPACE with quantum circuits}}

$L$ arbitrary language in PSPACE 

\hspace{2cm}{\Large $\Downarrow$}

family $(V_l)_{l\in\N}$ of quantum circuits such that

\quad $V_l$ decides for all input $x$ of length $l$ if $x\in L$:
\[
V_l^{r_l} (\ket{x} \otimes \ket{y} \otimes \ket{00\dots 0}) =
\ket{x} \otimes \ket{y \oplus f(x)\, {\rm mod}\, 2^{c_l}} 
\otimes \ket{00\dots 0}\,,
\]

\quad properties of $V_l$:

\quad\quad $V_l$ permutes only computational basis states

\quad\quad powers of $V_l$ produce orthogonal states (until recurrence)
\begin{eqnarray*}
f(x)=0 &
\ket{\Psi_x}\,,V_l\ket{\Psi_x}\,,\ldots,
V_l^{r_l-1}\ket{\Psi_x}\,,V_l^{r_l}\ket{\Psi_x}=\ket{\Psi_x} \\
f(x)=1 & 
{}\,\,\,\,\,\,\,\,\,\,\,\ket{\Psi_x}\,,V_l\ket{\Psi_x}\,,\ldots,
V_l^{{\red 2^{c_l}} r_l-1}\ket{\Psi_x}\,,V_l^{{\red 2^{c_l}} r_l}\ket{\Psi_x}=\ket{\Psi_x}
\end{eqnarray*}
\quad\quad{\red {\large $\Longrightarrow$ orbit length depends on $x\in L$}}
\end{slide}

\begin{slide}
{\bf {\large Local Hamiltonians/observables encoding the circuits}}

\vspace{0.25cm}
$T_k$: elementary quantum gates (Toffoli gates) of $V$ ($k=0,\ldots,s-1$)

\texttt{clock}: quantum register on $s$ qubits
\[
{\green \rightarrow}
\ket{1000\ldots 0}\,
{\red \mapsto}\,
\ket{0100\ldots0}\,
{\blue \mapsto}\,
\ket{0010\ldots 0}\,
\mapsto\,
\ldots\,\mapsto\,
\ket{000\ldots 01}\,
{\green \vdash}
\]

forward-time operator
\begin{eqnarray*}
F & = &
T_0\quad \otimes 
{\red \ket{1}_{2}\bra{0}_{1} \otimes \ket{0}_{1}\bra{1}_{0}}
\, +\, 
T_1 \otimes 
{\blue \ket{1}_{3}\bra{0}_{2} \otimes \ket{0}_{2}\bra{1}_{1}}
\, + \,\cdots\,+ \\
& &
T_{s-1} \otimes 
{\green \ket{1}_{1}\bra{0}_{0} \otimes \ket{0}_{s}\bra{1}_{s-1}}
\end{eqnarray*}

{\bf local Hamiltonian/observable}
\[
A:=F+F^\dagger
\]

\end{slide}

\begin{slide}
{\bf {\large Orbit properties}}

\begin{itemize}
\item $F$ acts as a cyclic shift on ${\cal
O}:=\big<F^j\ket{\Psi_x}\,\,|\,\, j\in\N\big>$
\item dimension of ${\cal O}$ is $d:=2^{{\red c \, f(x)}}\, r\, s $\quad
{\red depends on the solution}

$\Longrightarrow$ \quad eigenvalues and eigenvectors of $F$ restricted to
${\cal O}$ are
\[
e^{2\pi k/d}\quad\quad\ket{k}\quad\quad\mbox{ for } k=0,\ldots d-1
\]
$\Longrightarrow$ \quad $\ket{\Psi_x}$ uniform superposition of all
eigenvectors $\ket{k}$
\[
\ket{\Psi_x}:=\frac{1}{\sqrt{d}}\sum_{k=0}^{d-1}\ket{k}
\]
\item $F$ and $F^{\dagger}$ commute on ${\cal O}$ 

$\Longrightarrow$ \quad eigenvalues and eigenvectors of $A$ restricted
to ${\cal O}$ are
\[
2\cos(2\pi k/d)\quad\quad \ket{k}\quad\quad\mbox{ for } k=0,\ldots d-1
\]
\quad\quad $2$\quad EVs with multiplicity $1$ ($\equiv$ real EVs of $F$)

\quad $d-2$ EVs with multiplicity $2$ ($\equiv$ complex EVs of $F$)
\end{itemize}
\end{slide}

\begin{slide}
{\large {\bf Entropy of time-average}}

initial state:
\[
\ket{\Psi_x}\bra{\Psi_x}=\frac{1}{d}
\sum_{k,l} \ket{k}\bra{l}
\]
averaging over time destroys coherence between eigenspaces
\[
\bar{\rho}:=\frac{1}{d}\sum_{k,l}
[\cos(2\pi k/d)=\cos(2\pi l/d)]\, \ket{k}\bra{l}=
{\tiny \left(
\begin{array}{ccccccc}
1 &   &   &   &        &   &   \\
  & 1 &   &   &        &   &   \\
  &   & 1 & 1 &        &   &   \\
  &   & 1 & 1 &        &   &   \\
  &   &   &   & \ddots &   &   \\
  &   &   &   &        & 1 & 1 \\
  &   &   &   &        & 1 & 1 \\
\end{array}
\right)}
\]
\quad\quad entropy $S(\bar{\rho})=\log_2 d - (d-2)/d$

\vspace{0.25cm}
$d$ depends on whether $x$ is a ``yes'' or ``no'' instance

{\red $\Longrightarrow$ difference in entropy between ``yes'' and
``no'' instance at least $c/2$ bits}

\quad\quad $c$ size of the output register
\end{slide}

\begin{slide}
{\large {\bf Measurement of $A$}}

output register consists of one qubit

\vspace{0.25cm}
prepare $\ket{\Psi_x}$

perfect measurement of $A$ yields the results
\begin{itemize}
\item $2\cos(2\pi k / (2 r s))$,\quad $k=0,\ldots 2 rs -1$, if $x$
``yes'' instance
\item $2\cos(2\pi k / (\,\,r s))$,\quad $k=0,\ldots \,\,\, rs -1$, if $x$
``no'' instance
\end{itemize}
with equal probability

{\red $\Longrightarrow$ we can distinguish between both cases after few
samples}

\vspace{0.25cm}
imperfect measurement

{\red accuracy in the order $1/(rs)$ is sufficient to distinguish between
both cases}
\end{slide}

\end{document}

